\section[Campos sobre el retículo de incidencia]{Generación y localidad de los campos sobre el retículo de incidencia}
\label{sec:df-campos-reticulares}

La construcción reticular se realiza después de la selección del registro
dodecafásico y de sus incidencias. El origen marcado, el producto entero y
la forma de paridad que utiliza el espacio de estados proceden de esa
construcción. A continuación se explicitan las operaciones que producen
los campos, sus coeficientes y sus relaciones de localidad. El mismo
retículo interviene en la generación de estados, en los productos sobre
el vacío y en los conmutadores de todos los modos enteros.

\subsection{Residencia del origen marcado y persistencia de la incidencia}

% Fuentes: lean/incidencia/SelectedRegionalIncidence.lean;
% deltas/union/ConcreteNativeFourPlusOne.lean y SelectedExceptionalChain.lean.
% El antecedente SelectedFromRegionalInputs conserva en su propietario la
% interfaz S8 no ordenada y Lean.ofReduceBool del selector terminal. Este
% capítulo desarrolla sus salidas; no declara eliminadas esas dependencias.

Escribamos \(K_i\), \(0\leq i<12\), para las coordenadas del registro
seleccionado. Las dos lecturas de soporte son
\begin{equation}
 T_K=\{i:K_i\geq729\},\qquad
 H_K=\{i:p_i+e_i-f_i-K_i<0\},
 \label{eq:df-soportes-K}
\end{equation}
donde \(p_i,e_i,f_i\) son los bloques regionales utilizados por el lector
de preacarreo. El soporte superior utiliza el umbral entero \(729\);
el soporte negativo utiliza el signo del preacarreo antes de reducirlo.
Se conservan así dos operaciones diferentes sobre el mismo registro.

La selección ya establecida da
\[
 K=(234,543,140,729,659,824,621,58,914,794,146,601),
\]
como evaluación del registro, y produce la bandera
\(T_K\subset H_K\), con \(|T_K|=4\), \(|H_K|=6\).
Las incidencias regionales de propagación y autoescala, junto con el
soporte APP, determinan el conjunto
\[
 O_K=(E_{\mathrm{inc}}\cap F_{\mathrm{inc}})
       \setminus(H_K\cup A_{\mathrm{inc}}).
\]
La construcción de incidencia demuestra que \(O_K\) es un singleton.
Su único elemento, \(o\), fija el origen marcado que se conserva en
todo lo que sigue.

En la realización concreta de Golay, el soporte \(T_K\) corresponde a
\(\{4,6,8,14\}\). Las cuatro hexadas incidentes a ese soporte tienen
cuatro elevaciones octádicas; la estrella octádica completa tiene cinco
elementos. Su diferencia es una única octada adicional. El soporte
\(H_K\), por su parte, tiene una elevación octádica única en la misma
realización. Estas propiedades proceden de la realización residual
concreta y mantienen el origen regional del soporte seleccionado.

El vecino reticular marcado \(\Lambda_o\) obtenido de esa incidencia
tiene rango \(24\), forma bilineal entera y par, autodualidad integral
y mínimo cuadrático \(4\). El vector radial utilizado por la construcción
marcada tiene cuadrado \(54\). Denotaremos por
\(\langle x,y\rangle\in\mathbb Z\) la forma entera de \(\Lambda_o\).
Los teoremas de esta sección utilizan precisamente ese retículo y esa
forma ya construidos.

\subsection{Cociclo de paridad y álgebra reticular torcida}
\label{subsec:df-cociclo}

% Fuentes: deltas/voa/TriangularSignCocycle.lean,
% WittLatticeCocycle.lean y WittComplexAlgebra.lean.

Elijamos una base integral \((b_i)_{i=1}^{r}\) de \(\Lambda_o\), con
\(r=24\). Esta elección presenta las coordenadas de un retículo ya
obtenido. Su matriz de Gram es
\(G_{ij}=\langle b_i,b_j\rangle\). Para
\(x=\sum_i x_i b_i\), \(y=\sum_i y_i b_i\), definimos
\begin{equation}
 \tau(x,y)=\sum_{i<j}x_iG_{ij}y_j\pmod2,
 \qquad \varepsilon(x,y)=(-1)^{\tau(x,y)}.
 \label{eq:df-cociclo}
\end{equation}

\begin{proposition}[Identidad cocíclica y conmutador de paridad]
\label{prop:df-cociclo}
Se cumplen
\begin{align}
 \varepsilon(x,y)\varepsilon(x+y,z)
 &=\varepsilon(y,z)\varepsilon(x,y+z),\label{eq:df-identidad-cociclica}\\
 \varepsilon(x,y)\varepsilon(y,x)
 &=(-1)^{\langle x,y\rangle},\label{eq:df-conmutador-cociclo}\\
 \varepsilon(0,x)&=\varepsilon(x,0)=1.
\end{align}
\end{proposition}

\begin{proof}
La bilinealidad de \(\tau\) en \(\mathbb F_2\) da
\(\tau(x,y)+\tau(x+y,z)=\tau(y,z)+\tau(x,y+z)\), y al aplicar
\(a\mapsto(-1)^a\) se obtiene \eqref{eq:df-identidad-cociclica}.
La simetría de \(G\), junto con el carácter par de su diagonal, permite
reunir las dos sumas triangulares:
\[
 \tau(x,y)+\tau(y,x)
 =\sum_{i,j}x_iG_{ij}y_j
 =\langle x,y\rangle\pmod2.
\]
Esto prueba \eqref{eq:df-conmutador-cociclo}. Las identidades con cero
resultan directamente de \eqref{eq:df-cociclo}.
\end{proof}

El álgebra compleja torcida \(\mathbb C_{\varepsilon}[\Lambda_o]\)
consiste en las combinaciones lineales finitas de símbolos
\(\mathbf e^x\), con producto
\begin{equation}
 \mathbf e^x\mathbf e^y=\varepsilon(x,y)\mathbf e^{x+y}.
 \label{eq:df-producto-torcido}
\end{equation}
La identidad cocíclica demuestra la asociatividad y
\(\mathbf e^0\) es la unidad. La escritura \(\mathbf e^x\) designa el
elemento básico de carga \(x\), y conserva la etiqueta reticular completa.

\subsection{Espacio oscilatorio y representación de Heisenberg}
\label{subsec:df-osciladores}

% Fuente: deltas/voa/LatticeOscillatorFock.lean.

Sea \(\mathfrak h=\mathbb C\otimes_{\mathbb Z}\Lambda_o\). El espacio
oscilatorio algebraico es
\[
 \mathfrak h_-=
 \bigoplus_{n\geq1}\mathfrak h\,t^{-n},
 \qquad
 M(1)=\operatorname{Sym}_{\mathbb C}(\mathfrak h_-).
\]
El soporte de cada vector es finito, mientras que el conjunto de frecuencias
disponibles comprende todos los enteros positivos. El espacio de estados
que conserva simultáneamente ocupación oscilatoria y carga reticular es
\begin{equation}
 V_{\Lambda_o}=M(1)\otimes_{\mathbb C}
                  \mathbb C_{\varepsilon}[\Lambda_o],
 \qquad \mathbf1=1\otimes\mathbf e^0.
 \label{eq:df-espacio-estados}
\end{equation}

Para \(x\in\Lambda_o\) y \(n\geq1\), el operador \(x(-n)\) multiplica
por el generador \(x\,t^{-n}\). El operador \(x(n)\) es la derivación
de \(M(1)\) determinada por
\begin{equation}
 x(n)(y\,t^{-m})=n\,\delta_{n,m}\langle x,y\rangle.
 \label{eq:df-derivacion-oscilatoria}
\end{equation}
Ambos actúan como la identidad sobre el segundo factor de
\eqref{eq:df-espacio-estados}. El modo cero actúa sobre la carga mediante
\[
 x(0)(v\otimes\mathbf e^y)
 =\langle x,y\rangle v\otimes\mathbf e^y.
\]

\begin{theorem}[Relaciones de conmutación y no vacuidad]
\label{thm:df-heisenberg}
Para \(m,n\geq1\),
\begin{align}
 [x(n),y(-m)]&=n\delta_{n,m}\langle x,y\rangle I,
 \label{eq:df-CCR}\\
 [x(-n),y(-m)]&=0,
 & [x(n),y(m)]&=0.
\end{align}
Además, \(x(n)\mathbf1=0\) y \(\mathbf1\ne0\).
\end{theorem}

\begin{proof}
Si \(D\) es una derivación y \(M_a\) la multiplicación por \(a\), la
regla de Leibniz da \([D,M_a](v)=D(a)v\). Aplicada a
\eqref{eq:df-derivacion-oscilatoria}, produce \eqref{eq:df-CCR}.
Los multiplicadores conmutan porque el álgebra simétrica es conmutativa;
las derivaciones de coordenadas conmutan en los generadores y, por
Leibniz, en todos los monomios. La anulación del vacío se obtiene al
derivar la unidad. Finalmente, el funcional que toma el coeficiente del
monomio vacío y de la carga cero vale \(1\) sobre \(\mathbf1\), de modo
que el vacío es distinto de cero. La extensión por producto tensorial
conserva todas las identidades.
\end{proof}

Cada monomio oscilatorio tiene un peso entero
\(\operatorname{wt}(a)=\sum_{n,i}n a_{n,i}\), donde
\(a_{n,i}\) son sus ocupaciones. Esta graduación proporciona los límites
finitos de las sumas que aparecen en los coeficientes de campo.

\subsection{Coeficientes de creación y aniquilación en todos los grados}
\label{subsec:df-exponenciales}

% Fuentes: deltas/campos/LatticeCreationExponential.lean,
% LatticeAnnihilationExponential.lean y LatticeAnnihilationEnergy.lean.

Se construyen las dos series formales de operadores
\begin{equation}
 E^-_x(z)=\exp\!\left(\sum_{n\geq1}\frac{x(-n)}n z^n\right)
         =\sum_{d\geq0}C_x(d)z^d,
 \label{eq:df-creacion}
\end{equation}
\begin{equation}
 E^+_x(s)=\exp\!\left(-\sum_{n\geq1}\frac{x(n)}n s^n\right)
         =\sum_{r\geq0}A_x(r)s^r.
 \label{eq:df-aniquilacion}
\end{equation}
Estas exponenciales son operaciones formales sobre los modos ya
construidos. Cada coeficiente se define por una suma finita. En particular,
si \(P_x(z)=\sum_{n\geq1}x(-n)z^n/n\), entonces
\begin{equation}
 C_x(d)=\sum_{q=0}^{d}\frac1{q!}[z^d]P_x(z)^q,
 \label{eq:df-coeficiente-creacion}
\end{equation}
y se obtiene la fórmula análoga para \(A_x(r)\) sustituyendo el potencial
de aniquilación. El orden de composición de los endomorfismos se conserva
en cada potencia.

\begin{lemma}[Estabilización de los coeficientes]
\label{lem:df-estabilizacion}
Para cualquier potencial formal \(P\) con término constante cero,
\([z^d]P^q=0\) cuando \(q>d\). Por tanto, ampliar el límite superior
de \eqref{eq:df-coeficiente-creacion} a cualquier \(N\geq d\) conserva
su valor. Además,
\[
 C_x(0)=A_x(0)=I,
 \qquad A_x(r)1=0\quad(r>0).
\]
Si \(v\) tiene peso a lo sumo \(N\), entonces \(A_x(r)v=0\) para
\(r>N\).
\end{lemma}

\begin{proof}
Cada factor de \(P\) tiene grado al menos uno, por lo que un producto de
\(q\) factores tiene grado al menos \(q\). Este argumento utiliza
únicamente la graduación y se aplica también a coeficientes no
conmutativos. Las identidades en grado cero provienen de la potencia
vacía. Toda composición no vacía de modos de aniquilación anula la
unidad. Finalmente, un término de grado \(r\) en el potencial de
aniquilación reduce el peso oscilatorio en \(r\). Una reducción mayor
que el peso inicial da cero; por linealidad, la conclusión se conserva
en todo el subespacio de peso a lo sumo \(N\).
\end{proof}

Los primeros coeficientes ilustran la normalización:
\begin{align*}
 C_x(1)&=x(-1),&
 C_x(2)&=\tfrac12\bigl(x(-1)^2+x(-2)\bigr),\\
 A_x(1)&=-x(1),&
 A_x(2)&=\tfrac12\bigl(x(1)^2-x(2)\bigr).
\end{align*}
Cada denominador procede de la frecuencia o del orden de la potencia
formal. La forma de Gram que entra en los conmutadores sigue siendo la
forma entera del retículo marcado.

\subsection{Campo de carga y cota de truncamiento de Laurent}
\label{subsec:df-campo-carga}

% Fuente: deltas/campos/LatticeChargedVertexField.lean.

El campo asociado a una carga \(x\in\Lambda_o\) se escribe
\begin{equation}
 Y_x(z)=E^-_x(z)E^+_x(z^{-1})\,e_x z^{x(0)},
 \qquad Y_x(z)=\sum_{k\in\mathbb Z}Y_x[k]z^k,
 \label{eq:df-campo}
\end{equation}
donde \(e_x\mathbf e^y=\varepsilon(x,y)\mathbf e^{x+y}\).
La siguiente fórmula define sus coeficientes sobre una base monomial y
explicita la finitud de cada operación.

\begin{definition}[Coeficiente de campo sobre un estado cargado]
Si \(v\) es un monomio de peso \(N\), se define
\begin{equation}
 Y_x[k](v\otimes\mathbf e^y)
 =\varepsilon(x,y)
 \sum_{j=0}^{N}
 C_x\bigl(k-\langle x,y\rangle+j\bigr)A_x(j)v
       \otimes\mathbf e^{x+y},
 \label{eq:df-coeficiente-campo}
\end{equation}
con la convención \(C_x(d)=0\) para \(d<0\). Se extiende linealmente
a \(V_{\Lambda_o}\).
\end{definition}

\begin{theorem}[Estabilidad, truncamiento y creación de estados de carga]
\label{thm:df-campo-truncamiento}
La fórmula \eqref{eq:df-coeficiente-campo} es independiente de cualquier
cota de peso superior a \(N\). Se cumple
\begin{equation}
 k<\langle x,y\rangle-N
 \quad\Longrightarrow\quad
 Y_x[k](v\otimes\mathbf e^y)=0.
 \label{eq:df-truncamiento-laurent}
\end{equation}
Cada estado admite, por tanto, una cota inferior de Laurent. Sobre el
vacío,
\begin{equation}
 Y_x[k]\mathbf1=
 \begin{cases}
 C_x(k)1\otimes\mathbf e^x,&k\geq0,\\
 0,&k<0.
 \end{cases}
 \label{eq:df-campo-vacio}
\end{equation}
En particular, \(Y_x[0]\mathbf1=1\otimes\mathbf e^x\).
\end{theorem}

\begin{proof}
La ampliación de la cota añade términos \(A_x(j)v\) con \(j>N\), que
son cero por el lema~\ref{lem:df-estabilizacion}. Si se cumple la
desigualdad de \eqref{eq:df-truncamiento-laurent}, todos los índices
\(k-\langle x,y\rangle+j\), con \(j\leq N\), son negativos; cada
sumando se anula. Una combinación lineal finita de monomios conserva
la menor de sus cotas, lo que prueba la afirmación para cualquier
estado. Sobre \(\mathbf1\), sólo subsiste \(j=0\), y
\(\varepsilon(x,0)=1\), \(\langle x,0\rangle=0\). Se obtiene
\eqref{eq:df-campo-vacio}.
\end{proof}

\begin{theorem}[Generación de todo el espacio reticular de estados]
\label{thm:df-generacion-campos}
El subespacio generado a partir del vacío por composiciones finitas de
los operadores \(Y_x[0]\), \(x\in\Lambda_o\), y de todos los modos
de creación \(b_i(-n)\), \(n\geq1\), es \(V_{\Lambda_o}\).
\end{theorem}

\begin{proof}
Por \eqref{eq:df-campo-vacio}, \(Y_x[0]\) aplicado al vacío produce
el estado fundamental de cualquier carga \(x\). Los operadores de
creación multiplican por los generadores del álgebra simétrica;
composiciones finitas de ellos producen todos los monomios sobre ese
estado fundamental. Los monomios oscilatorios por los elementos
\(\mathbf e^x\) constituyen una base algebraica de
\eqref{eq:df-espacio-estados}. Su envolvente lineal es, por tanto, todo
el espacio. El argumento comprende cualquier frecuencia, ocupación
finita y carga reticular.
\end{proof}

\begin{corollary}[Determinación de los entrelazadores por el vacío]
Dos endomorfismos de \(V_{\Lambda_o}\) que conmutan con todos los
generadores del teorema anterior y coinciden sobre el vacío son iguales.
\end{corollary}

\begin{proof}
La conmutación transporta la igualdad en el vacío a cualquier
composición finita de generadores. La generación total y la linealidad
extienden esa igualdad a todo el espacio.
\end{proof}

\subsection{Contracción escalar y ordenación normal de los productos}
\label{subsec:df-contraccion}

% Fuentes: deltas/generacion/LatticeExponentialContraction.lean,
% LatticeContractionFactor.lean, LatticeNormalOrdering.lean;
% deltas/productos/LatticeVacuumContraction.lean.

El producto entero \(p=\langle x,y\rangle\) fija una sucesión escalar
por la recurrencia
\begin{equation}
 s_p(0)=1,\qquad
 (r+1)s_p(r+1)=(r-p)s_p(r).
 \label{eq:df-contraccion-recurrencia}
\end{equation}
La recurrencia determina cada coeficiente a partir del anterior. Su
expresión cerrada es
\(s_p(r)=(-1)^r\binom pr\), con el binomial entero generalizado.
Los primeros términos son
\[
 1,\quad -p,\quad \frac{p(p-1)}2,\quad
 -\frac{p(p-1)(p-2)}6.
\]

\begin{proposition}[Factor formal de contracción]
\label{prop:df-factor-contraccion}
La serie \(S_p(u)=\sum_{r\geq0}s_p(r)u^r\) satisface
\begin{equation}
 (1-u)S_p'(u)=-pS_p(u),\qquad S_p(0)=1.
 \label{eq:df-ecuacion-factor}
\end{equation}
Es la única solución formal de esas condiciones. Además,
\[
 S_{p+q}=S_pS_q,qquad S_0=1,qquad S_1=1-u,qquad
 S_pS_{-p}=1.
\]
Para \(p\geq0\), \(S_p(u)=(1-u)^p\), y sus coeficientes de grado
superior a \(p\) se anulan.
\end{proposition}

\begin{proof}
Comparar el coeficiente de \(u^r\) en
\eqref{eq:df-ecuacion-factor} reproduce exactamente
\eqref{eq:df-contraccion-recurrencia}. Como \(r+1\ne0\), los
coeficientes quedan determinados inductivamente. La regla de Leibniz
muestra que \(S_pS_q\) satisface la misma ecuación que \(S_{p+q}\)
y tiene el mismo término constante; la unicidad da la identidad
multiplicativa. Los casos \(0,1\) se verifican directamente. La
identidad inversa resulta del caso \(q=-p\), y las potencias naturales
se obtienen por inducción. Si \(p\) es natural, el paso \(r=p\) de
la recurrencia produce cero y los pasos posteriores conservan ese cero.
\end{proof}

\begin{theorem}[Ordenación normal en todos los grados]
\label{thm:df-ordenacion}
Las series construidas a partir de los modos satisfacen
\begin{equation}
 E^+_x(s)E^-_y(z)
 =S_{\langle x,y\rangle}(sz),E^-_y(z)E^+_x(s).
 \label{eq:df-ordenacion-normal}
\end{equation}
En coeficientes, para \(r,d\geq0\),
\begin{equation}
 A_x(r)C_y(d)
 =\sum_{i=0}^{\min(r,d)}
 s_{\langle x,y\rangle}(i),
 C_y(d-i)A_x(r-i).
 \label{eq:df-ordenacion-coeficientes}
\end{equation}
\end{theorem}

\begin{proof}
Sean
\(Q_x(s)=-\sum_{n\geq1}x(n)s^n/n\) y
\(P_y(z)=\sum_{m\geq1}y(-m)z^m/m\).
Por \eqref{eq:df-CCR},
\[
 [Q_x(s),P_y(z)]
 =-\langle x,y\rangle\sum_{n\geq1}\frac{(sz)^n}{n}\,I.
\]
El conmutador es central. La identidad exponencial para dos potenciales
con conmutador central se demuestra aquí coeficiente a coeficiente:
conmutar una potencia de \(Q_x\) a través de una potencia de \(P_y\)
añade los términos escalares de ese conmutador; las sumas a cada
bigrado son finitas por el lema~\ref{lem:df-estabilizacion}. El factor
escalar resultante tiene término constante uno y satisface
\eqref{eq:df-ecuacion-factor}, con
\(p=\langle x,y\rangle\); por la unicidad de la
proposición~\ref{prop:df-factor-contraccion}, es \(S_p(sz)\).
Esto demuestra \eqref{eq:df-ordenacion-normal}. La extracción del
coeficiente \(s^rz^d\) da
\eqref{eq:df-ordenacion-coeficientes}; el factor de contracción consume
el mismo grado \(i\) en ambas variables, de donde procede el límite
\(i\leq\min(r,d)\).
\end{proof}

\begin{theorem}[Contracción exacta de los estados creados desde el vacío]
\label{thm:df-contraccion-vacio}
Para toda pareja de cargas \(x,y\) y todos los grados \(r,d\),
\begin{equation}
 A_x(r)C_y(d)1=
 \begin{cases}
 s_{\langle x,y\rangle}(r),C_y(d-r)1,&r\leq d,\\
 0,&r>d.
 \end{cases}
 \label{eq:df-contraccion-vacio}
\end{equation}
En la diagonal \(r=d\), el resultado es
\(s_{\langle x,y\rangle}(r)1\). Si
\(\langle x,y\rangle=p\geq0\), también se anula para \(r>p\).
\end{theorem}

\begin{proof}
Evaluamos \eqref{eq:df-ordenacion-coeficientes} sobre \(1\).
El lema~\ref{lem:df-estabilizacion} anula todos los sumandos con
\(r-i>0\). Subsiste únicamente \(i=r\), siempre que \(r\leq d\).
Para \(r=d\), \(C_y(0)=I\). La última anulación sigue de la cota
polinómica de \(S_p\).
\end{proof}

Por ejemplo, \(A_x(1)C_y(1)1=-\langle x,y\rangle1\).
Para \(p=-1\), la recurrencia da \(s_{-1}(r)=1\) en todos los grados;
para \(p=2\), da \((1,-2,1,0,\ldots)\). Estas dos evaluaciones
exhiben, respectivamente, una contracción de orden ilimitado y una
cancelación polinómica finita, ambas determinadas por la misma forma
entera del retículo.

\subsection{Conmutador mixto en todos los modos enteros}
\label{subsec:df-conmutador-mixto}

% Fuentes: deltas/integracion/LatticePositiveMixedFields.lean,
% LatticeNegativeMixedFields.lean y LatticeAllMixedFields.lean.

Para \(x\in\Lambda_o\), escribimos \(H_x(m)=x(m)\), con
\(m\in\mathbb Z\), incluyendo los modos negativos de creación,
el modo cero de carga y los positivos de aniquilación.

\begin{theorem}[Conmutador mixto integral]
\label{thm:df-conmutador-mixto}
Para todas las cargas \(x,y\) y todos los enteros \(m,k\),
\begin{equation}
 [H_x(m),Y_y[k]]
 =\langle x,y\rangle,Y_y[k-m].
 \label{eq:df-conmutador-mixto}
\end{equation}
\end{theorem}

\begin{proof}
Tratamos separadamente los tres signos de \(m\), conservando en cada
caso los coeficientes de campo de \eqref{eq:df-coeficiente-campo}.

Si \(m>0\), la relación de Heisenberg da
\[
 [x(m),C_y(d)]=
 \begin{cases}
 \langle x,y\rangle C_y(d-m),&d\geq m,\\
 0,&d<m.
 \end{cases}
\]
Los modos positivos conmutan con los coeficientes de aniquilación y
con el desplazamiento de carga. Al conmutar término a término en
\eqref{eq:df-coeficiente-campo}, el grado de creación disminuye en
\(m\), que equivale a reemplazar \(k\) por \(k-m\). El corte de peso
permanece válido porque la aniquilación disminuye el peso.

Si \(m=0\), el campo desplaza la carga \(z\) a \(y+z\). La diferencia
de los autovalores del modo cero es
\(\langle x,y+z\rangle-\langle x,z\rangle=\langle x,y\rangle\).
Los operadores oscilatorios preservan la carga, de modo que se obtiene
\eqref{eq:df-conmutador-mixto} con \(k-m=k\).

Finalmente, sea \(m=-a\), con \(a>0\). Los modos de creación conmutan
entre sí, mientras que la relación derivada de
\eqref{eq:df-CCR} y \eqref{eq:df-aniquilacion} es
\[
 [x(-a),A_y(j)]=
 \begin{cases}
 \langle x,y\rangle A_y(j-a),&j\geq a,\\
 0,&j<a.
 \end{cases}
\]
Si el estado inicial tiene peso a lo sumo \(N\), después de crear el
modo \(-a\) tiene peso a lo sumo \(N+a\). Usamos este corte común en
los dos términos del conmutador. Los sumandos con \(j<a\) se anulan;
reindexamos los restantes mediante \(j=r+a\), \(0\leq r\leq N\).
El grado de creación se transforma en
\(k-\langle y,z\rangle+r+a\), exactamente el coeficiente de campo
con índice \(k+a=k-m\). La identidad queda probada sobre los estados
monomiales cargados y se extiende por linealidad a todo el espacio.
\end{proof}

El desplazamiento del índice de campo distingue correctamente la
orientación de los modos: un modo positivo resta su frecuencia, un modo
negativo la suma y el modo cero conserva el índice. Los tres casos son
realizaciones de una única identidad en \(\mathbb Z\).

\subsection{Localidad mixta por cancelación coeficiente a coeficiente}
\label{subsec:df-localidad-mixta}

% Fuentes: deltas/integracion/LatticeMixedLocalityCriterion.lean,
% LatticeMixedLocality.lean y SelectedMixedFieldLocality.lean.

Consideremos el campo de Heisenberg
\(H_x(z)=\sum_{m\in\mathbb Z}H_x(m)z^{-m-1}\).
Para una serie bivariada de coeficientes \(F(a,b)\), la multiplicación
por \(z-w\) actúa como
\begin{equation}
 (\mathcal C F)(a,b)=F(a-1,b)-F(a,b-1).
 \label{eq:df-operador-cruce}
\end{equation}
Las composiciones ordenadas de los dos campos tienen coeficientes
\[
 F(a,b)=H_x(-a-1)Y_y[b],\qquad
 G(a,b)=Y_y[b]H_x(-a-1).
\]

\begin{theorem}[Localidad mixta de orden uno]
\label{thm:df-localidad-mixta}
Los campos construidos satisfacen
\begin{equation}
 (z-w)H_x(z)Y_y(w)=(z-w)Y_y(w)H_x(z).
 \label{eq:df-localidad-mixta}
\end{equation}
La igualdad es una identidad de coeficientes endomorfos en todo
\(V_{\Lambda_o}\).
\end{theorem}

\begin{proof}
Por \eqref{eq:df-operador-cruce}, el coeficiente \((a,b)\) de la
diferencia es
\[
 [H_x(-a),Y_y[b]]-[H_x(-a-1),Y_y[b-1]].
\]
El teorema~\ref{thm:df-conmutador-mixto} convierte ambos términos en
\(\langle x,y\rangle Y_y[a+b]\). Su resta es cero para todos los
enteros \(a,b\). La identidad se cumple en los endomorfismos, por lo
que vale simultáneamente para todos los estados.
\end{proof}

\subsection{Localidad de los campos cargados y conservación del origen}
\label{subsec:df-localidad-cargada}

% Fuentes: deltas/localidad/LatticeFactorConvolution.lean,
% LatticeChargedLocality.lean, LatticeChargedLocalityFull.lean,
% SelectedFieldLocality.lean y deltas/integracion/SelectedMixedFieldLocality.lean.

La ordenación normal permite demostrar también la localidad de dos
campos de carga. Sea \(p=\langle x,y\rangle\) y
\(N_p=\max(0,-p)\). El factor de contracción de los dos órdenes se
expande en las regiones formales \(w/z\) y \(z/w\); su paridad relativa
está fijada por \eqref{eq:df-conmutador-cociclo}.

\begin{theorem}[Localidad uniforme de los campos cargados]
\label{thm:df-localidad-cargada}
Para cualesquiera \(x,y\in\Lambda_o\),
\begin{equation}
 (z-w)^{N_p}Y_x(z)Y_y(w)
 =(z-w)^{N_p}Y_y(w)Y_x(z),
 \qquad N_p=\max(0,-\langle x,y\rangle).
 \label{eq:df-localidad-cargada}
\end{equation}
El exponente depende de las cargas y es válido para todos los estados.
\end{theorem}

\begin{proof}
Sobre un estado de peso acotado, los productos de los coeficientes de
campo se expresan mediante \eqref{eq:df-ordenacion-coeficientes}, los
desplazamientos de carga y el cociclo. Las sumas por coeficiente son
finitas: la aniquilación tiene corte de peso y las dos variables del
producto ordenado tienen una cota inferior. La identidad cocíclica
reúne los desplazamientos en el mismo elemento de carga \(x+y\), y
la identidad de conmutador aporta el factor \((-1)^p\) entre los
dos órdenes.

Los factores escalares resultantes son las dos expansiones de
\((z-w)^p\). Si \(p\geq0\), ambas son el mismo polinomio y coinciden
sin multiplicación adicional. Si \(p<0\), multiplicar por
\((z-w)^{-p}\) transforma ambas expansiones en el factor uno.
Formalmente, cada aplicación de \(\mathcal C\) eleva en uno el
exponente de ambas expansiones; después de \(-p\) aplicaciones se
alcanza el exponente cero. Esta cancelación es independiente del corte
de peso elegido y de la carga del estado sobre el que actúan los
campos. Se obtiene así \eqref{eq:df-localidad-cargada} sobre todo
estado monomial cargado, y la linealidad la extiende a
\(V_{\Lambda_o}\).
\end{proof}

\begin{theorem}[Composición conservativa de incidencia, generación y localidad]
\label{thm:df-composicion-campos}
Fijado el origen \(o\) producido por la incidencia del registro
seleccionado, el mismo espacio \(V_{\Lambda_o}\) satisface conjuntamente:
la generación desde el vacío del
teorema~\ref{thm:df-generacion-campos}, los productos de todos los grados
del teorema~\ref{thm:df-contraccion-vacio}, la localidad cargada de
\eqref{eq:df-localidad-cargada} y la localidad mixta de
\eqref{eq:df-localidad-mixta}. Su forma entera, su cociclo, su vacío
y sus operadores de carga son los definidos a partir de
\(\Lambda_o\) en esta sección.
\end{theorem}

\begin{proof}
Todos los enunciados anteriores están cuantificados sobre el mismo
retículo y la misma forma entera. Al especializarlos al origen
\(o\) ya seleccionado, sus objetos coinciden por sus definiciones:
el cociclo usa la base de \(\Lambda_o\), los osciladores usan su
Gram, los campos utilizan esos osciladores y ese cociclo, y los
productos utilizan esos mismos coeficientes de campo. La conjunción
de los teoremas preserva, por tanto, la bandera de incidencia de
\eqref{eq:df-soportes-K} y todas las dependencias que determinaron
el origen. Ninguna especialización introduce una nueva elección
reticular.
\end{proof}

La concatenación demostrativa distingue con precisión sus operaciones:
los bloques regionales determinan prefijos, cilindros y firmas; la
selección dodecafásica y la incidencia producen el origen reticular;
el retículo determina la forma entera y el cociclo; éstos fijan los
modos, los campos y sus productos. La estructura de campos resultante
queda así provista de generación total y de identidades de localidad
en el espacio algebraico completo. Los capítulos dedicados a la
reconstrucción de la estructura de operadores de vértice, al sector
torcido y a la realización excepcional utilizan estos resultados con
sus respectivas hipótesis adicionales, manteniendo esta misma
procedencia generativa.
